\chapter{结论与建议}\label{ux7ed3ux8bbaux4e0eux5efaux8bae}

\section{主要研究结论}\label{ux4e3bux8981ux7814ux7a76ux7ed3ux8bba}

本研究以佛山市南海区为对象，构建了``文化---旅游''双谱系对照框架，并以文化记忆指数、官方认证指数、旅游热度指数和文化---旅游错位指数为主要工具，分析了 165 条文化载体和 500 m 网格尺度下的空间格局。需要说明的是，本文结论主要指向文旅融合潜力和错位关系识别，而不是对文旅融合成熟度的完整评价。主要结论如下。

\textbf{结论一：南海区文化资源与旅游产品之间存在较明显的错位。}

在载体尺度上，165 条文化载体中，一般耦合区 57 条，占 34.5\%；沉睡潜力区 54 条，占 32.7\%；空心景点区 49 条，占 29.7\%；核心耦合区仅 5 条，占 3.0\%。在网格尺度上，1 跳知识图谱口径识别出 170 个沉睡潜力网格、34 个空心景点网格和 8 个核心耦合网格。这说明南海区文旅融合的关键问题并不是文化资源或旅游产品单方面不足，而是二者在空间载体和体验转化上尚未形成充分对应。其中，核心耦合区应理解为文化基础和旅游热度同时较强的优先深化对象，而不是融合已经完成的成熟区。

\textbf{结论二：``东旅西文''的空间结构较为稳定。}

POI 高值区主要集中在桂城、大沥、狮山等东北部城镇化片区，而文化载体与文化厚度高值更多出现在西樵、九江、丹灶等西南部片区。镇街级相关结果也显示，物质载体数与 POI 总量呈中等负相关。由此可以看到，南海区并不是简单的``哪里文化多、哪里旅游就强''，而是存在较清楚的空间方向性错位。

\textbf{结论三：物质文化载体适合作为主桥梁，非遗更适合作为动态补充。}

物质文化载体通常具备较明确的空间位置和官方认定信息，便于与知识图谱、POI 和网格分析相连接，因此适合作为本研究的主操作化桥梁。非遗项目中有一部分能够依托场馆、训练基地、村落或节庆空间参与空间匹配，但也有相当一部分只能定位到传承人或镇街层面。对于后者，更适合在规划建议中以活动植入、展示转译和线路组织的方式发挥作用，而不宜强行视为固定景点。

\textbf{结论四：相关分析支持错位判断，但不能作因果解释。}

载体级相关矩阵显示，CMI 与 MI 呈中等负相关（r = −0.495），OAI 与 MI 呈中等负相关（r = −0.440），而 500 m 范围内 POI 数量与 CMI、OAI 的相关性接近于零。这一结果从量化层面支持了``文化高、旅游低''和``旅游高、文化支撑不足''并存的判断。与此同时，THI 由 POI、评分和评论量构成，主要反映旅游热度，不直接反映游客文化识别；镇街级样本只有 7 个，A 级景区样本只有 16 个，相关系数也只适合提示关系，不应被解释为严格因果。补充的权重敏感性检验显示，四组替代口径下分类一致比例均在 89.7\% 以上，说明主要空间判断具有一定稳定性；评论文化识别度抽样则说明，部分核心载体已经出现文化关键词，但仍不能把旅游热度直接等同于文化体验深度。

\textbf{结论五：西南廊道是潜力释放重点，东北都市圈更需要文化内容补强。}

1 跳知识图谱口径下，西樵、九江、丹灶三镇合计 147 个沉睡潜力网格，占全区 1 跳沉睡潜力网格的 86.5\%，说明西南片区具备集中释放文化潜力的空间基础。相对而言，桂城、大沥等东北部片区旅游热度较高，但文化厚度相对不足，更需要通过非遗活动、文化市集、解说系统和主题事件补足地方文化内容。

\textbf{结论六：九江镇需要同时考虑真实转化不足与数据口径低估。}

九江镇文化线索较密，但 POI 数量和平台评论热度偏低。这一结果一方面反映其交通联系、景区化程度和旅游产品组织仍有不足；另一方面也受到本文未纳入餐饮门店和餐饮评价平台数据的影响。对于九江这样的美食和民俗体验型片区，后续研究和规划判断需要把``数据低估''和``实质短板''区分开来。

\section{文旅融合潜力释放与提升建议}\label{ux6587ux65c5ux878dux5408ux6f5cux529bux91caux653eux4e0eux63d0ux5347ux5efaux8bae}

从风景园林和城乡规划的角度看，文旅融合潜力释放不能只停留在资源宣传层面，而需要落实到载体、游线、场地和运营机制上。后续干预也不宜只按``热度高低''排序，而应同时考虑文化基础、旅游热度、空间可达性和内容转化条件。本文据此提出以下建议。

表 7.1　文旅融合潜力释放梯队与规划方向

{\def\LTcaptype{none} % do not increment counter
\begin{longtable}[]{@{}
  >{\centering\arraybackslash}p{(\linewidth - 6\tabcolsep) * \real{0.3077}}
  >{\raggedright\arraybackslash}p{(\linewidth - 6\tabcolsep) * \real{0.2308}}
  >{\raggedright\arraybackslash}p{(\linewidth - 6\tabcolsep) * \real{0.2308}}
  >{\raggedright\arraybackslash}p{(\linewidth - 6\tabcolsep) * \real{0.2308}}@{}}
\toprule\noalign{}
\begin{minipage}[b]{\linewidth}\centering
梯队
\end{minipage} & \begin{minipage}[b]{\linewidth}\raggedright
代表载体
\end{minipage} & \begin{minipage}[b]{\linewidth}\raggedright
数据特征
\end{minipage} & \begin{minipage}[b]{\linewidth}\raggedright
规划方向
\end{minipage} \\
\midrule\noalign{}
\endhead
\bottomrule\noalign{}
\endlastfoot
第一梯队：优先深化型 & 云泉仙馆、西樵山百步云梯、九江吴家大院、西樵山抗日阵亡将士暨死难同胞纪念碑、平洲传统玉器制作技艺 & CMI、OAI 与 THI 均较高，MI 接近 0，属于核心耦合区 & 作为品牌节点、线路起讫点和深度体验节点，继续开展评论语义复核与产品深化 \\
第二梯队：条件改善型 & 西樵镇松塘村、九江镇烟桥烟南村、丹灶镇仙岗社区、赛龙舟（九江传统龙舟）、西樵山仙姑亭 & 文化记忆和官方认证较强，但 THI 偏低，MI 显著为负 & 优先补可达性、展示空间、公共服务、活动运营和线上可见度 \\
第三梯队：文化补强型 & 南海农谚、湖山胜迹门楼、小云亭、象林塔、光分亭 & 旅游热度较高，但文化记忆或官方支撑偏弱，MI 显著为正 & 补地方叙事、非遗活动、解说系统和可被游客识别的文化内容 \\
\end{longtable}
}

\subsection{载体层面：按错位类型分级干预}\label{ux8f7dux4f53ux5c42ux9762ux6309ux9519ux4f4dux7c7bux578bux5206ux7ea7ux5e72ux9884}

对于沉睡潜力区，应优先改善可达性、展示空间和解说系统。松塘村、烟桥古村、仙岗村等历史村落已经具备一定文化识别基础，下一步更需要完善游线入口、公共服务、导览标识和停留节点，使文化记忆能够转化为游客可以感知的游览过程。

对于空心景点区，应重点补足本土文化叙事。桂城、大沥等现代公园、商业休闲空间和城市开放空间已经具备较强人流基础，但地方文化内容不足。可以通过非遗快闪、周末文化市集、主题展陈和临时性公共艺术等方式，把文化内容嵌入日常休闲空间。

对于核心耦合区，应将其作为优先深化对象和线路组织基础节点，而不是简单认定为成熟示范区。云泉仙馆、西樵山百步云梯、九江吴家大院、平洲传统玉器制作技艺等已具备一定文化和旅游双重基础，适合进一步承担片区线路的起点、转换点或主题节点。下一步还应结合评论语义、现场解说和产品内容，判断游客是否真正识别到其地方文化价值。

\subsection{片区层面：形成西南廊道与东北都市圈的差异化策略}\label{ux7247ux533aux5c42ux9762ux5f62ux6210ux897fux5357ux5ecaux9053ux4e0eux4e1cux5317ux90fdux5e02ux5708ux7684ux5deeux5f02ux5316ux7b56ux7565}

西南廊道应以``文化资源转化''为重点。西樵---松塘---烟桥---九江---丹灶一线可围绕理学文化、古村落、龙舟、美食非遗和水乡景观组织复合游线，在烟桥、松塘、仙岗等 MI 负值较大的节点周边优先布设小型展示空间、活态传习点和研学停留点。

九江镇应突出``美食非遗 + 龙舟研学 + 古村游线''的组合。九江双蒸博物馆、吴家大院、烟桥古村、龙舟训练和水乡餐饮空间可以共同构成片区骨架。对于当前模型识别出的沉睡潜力网格，可优先补充文化故事卡、步行导览、餐饮体验挂接和节庆活动节点，同时在后续数据工作中补充餐饮评价和夜间活力数据。

东北都市圈应以``文化内容注入''为重点。桂城、大沥、狮山和里水已经拥有较强 POI 密度和日常消费基础，适合把非遗展示、地方历史叙事和公共活动嵌入公园、湖区、商圈和社区开放空间，使现代旅游供给不只停留在环境消费和休闲消费层面。

\subsection{管理层面：将错位指数纳入动态监测}\label{ux7ba1ux7406ux5c42ux9762ux5c06ux9519ux4f4dux6307ux6570ux7eb3ux5165ux52a8ux6001ux76d1ux6d4b}

建议将 CMI、OAI、THI 和 MI 作为南海区文旅融合潜力的动态监测指标，定期更新载体级和网格级结果。载体级结果可用于形成项目清单和优先级排序，网格级结果可用于识别连续片区、游线节点和公共服务短板。监测时应保持指标口径相对稳定，同时补充餐饮评论、客流热力、活动数据和评论文化关键词识别，以便观察干预措施是否真正改变了文化资源与旅游供给之间的对应关系。

\section{研究不足与展望}\label{ux7814ux7a76ux4e0dux8db3ux4e0eux5c55ux671b}

\subsection{数据层面}\label{ux6570ux636eux5c42ux9762}

本研究仍存在若干数据限制。第一，典籍 OCR 和实体抽取仍可能存在识别误差，部分古地名、人物名和官职名需要后续人工校核。第二，评论数据是单一时点截面，且未纳入大众点评、美团、百度地图等平台，对九江等餐饮和生活体验主导的片区可能存在低估。第三，本文已经补充评论文化识别度抽样，但该检验主要依靠 500 m 关联评论和关键词匹配，尚不是全量语义分类模型，因此仍不能直接评价文旅融合深度。第四，55 项主要依托传承人登记地的非遗项目空间边界较弱，尚不能像物质载体一样完整参与主桥梁匹配。第五，镇街边界和网格划分仍存在少量空间误差，部分网格被标记为未标注，但从整体结果看并未改变主要空间格局。

\subsection{方法层面}\label{ux65b9ux6cd5ux5c42ux9762}

方法上，错位指数中的 CMI 与 OAI 采用等权设置，THI 中 POI 数量、评分和评论量也采用预设权重。这种处理有利于保持本科论文阶段的透明性和可复核性，但仍然具有简化性。本文已补充四组权重敏感性检验，结果显示主要分类具有一定稳定性，但这仍不是全面的参数搜索，也没有引入专家打分、游客问卷或多元模型进行校准。对于处在阈值边界附近的样本，后续仍需结合现场调研和产品运营情况作进一步判断。同时，镇街级相关分析样本量仅为 7，A 级景区样本量仅为 16，相关结果只能作为提示，不宜作强因果解释。

\subsection{展望}\label{ux5c55ux671b}

后续研究可以从三个方向继续推进。第一，补充餐饮评价、夜间热度、节庆活动和客流数据，以更准确地识别九江等片区的真实旅游活力。第二，扩大评论文化识别度分析，从关键词抽样走向全量语义分类、情感判断和主题模型识别，用来判断``热度''是否真正转化为文化感知。第三，深化知识图谱的路径分析，从``主题---地点---关联实体''的角度筛选文旅 IP 和潜力项目组合，并将本研究框架拓展到更细的街区和地块尺度，在佛山其他区县或珠三角同尺度区县中进行比较，以检验方法的适用性和可迁移性。
